\section{Examples Belong to the Hardware They Need}
\label{sec:bsp-examples}

The most practical BSP rule in this book is about where example files live, because that is where a
layout either documents itself or lies.

\subsection{The Rule}

\begin{quote}
    An example lives in the crate whose \textbf{smallest} hardware requirement it satisfies.
\end{quote}

\begin{table}[htbp]
    \centering
    \footnotesize
    \begin{tabular}{@{}p{0.36\textwidth}p{0.28\textwidth}p{0.32\textwidth}@{}}
        \toprule
        \textbf{Example} & \textbf{Lives in} & \textbf{Needs} \\
        \midrule
        \codeinline{module\_led\_blink} & \codeinline{xiao-mg24-bsp} & the module \\
        \codeinline{module\_imu} & \codeinline{xiao-mg24-bsp} (\codeinline{sense}) & the module, Sense variant \\
        \codeinline{base\_env\_clock} & \codeinline{xiao-expansion-base-bsp} & module + carrier \\
        \codeinline{base\_fall\_detector} & \codeinline{xiao-expansion-base-bsp} & module + carrier (panel, buzzer) \\
        \bottomrule
    \end{tabular}
    \caption{The path of the file is the bill of materials.}
    \label{tab:bsp-examples}
\end{table}

\subsection{The Three Ways to Get It Wrong}

\begin{enumerate}
    \item \textbf{Everything in the carrier crate.} Convenient, and it makes the module crate's API
          unexercised: this project's module facade had no examples for weeks, and its documentation
          described a pin that did not exist.
    \item \textbf{Everything in the module crate.} Impossible without the carrier's devices, so the
          examples would need \codeinline{cfg} guards and would prove neither level.
    \item \textbf{A flat examples/ directory at the workspace root.} The path stops telling you
          anything, and \codeinline{cargo} stops checking that the crate still builds its own examples.
\end{enumerate}

\subsection{The Check That Keeps It True}

\begin{minted}{bash}
# every crate that exposes hardware must carry at least one example of its own level
for dir in crates/*/ libraries/*/; do
  [ -f "${dir}Cargo.toml" ] || continue
  n=$(find "$dir/examples" -maxdepth 1 -name '*.rs' | wc -l)
  [ "$n" -gt 0 ] || fail "$(basename "$dir"): no examples -- its API is not built by anything"
done
\end{minted}

That check is why this project's module facade has a \codeinline{module\_led\_blink} that uses only the
module, and why the carrier's examples are in the carrier. A rule enforced by the file system is a rule
nobody has to remember.

\subsection{The Half-Example That Taught the Rule}

A fall detector was written in the HAL's examples directory. It used the IMU, the panel and the buzzer:
two boards' worth of hardware, in the crate that knows neither. It compiled, because the HAL happened to
contain the wiring at the time — which was itself the defect the split fixed.

After the split, the same program had an obvious home: the carrier crate, which depends on the module
and therefore has both the IMU (through the module) and the panel (its own). Moving it changed no code.
That is the sign of a layout that was right: the program did not have to change to be correctly placed.
